\documentclass[11pt,a4paper]{article}
\usepackage[utf8]{inputenc}
\usepackage[vietnamese]{babel}
\usepackage{amsmath,amssymb}
\usepackage{geometry}
\geometry{a4paper, total={180mm,260mm}, left=15mm, right=15mm, top=18mm, bottom=18mm}
\usepackage{tikz}
\usetikzlibrary{shapes.geometric, arrows.meta, positioning, calc}
\usepackage{tcolorbox}
\tcbuselibrary{skins,breakable}
\usepackage{xcolor}
\usepackage{multicol}
\usepackage{array}
\usepackage{fontawesome}
\usepackage{fancyhdr}

\pagestyle{fancy}
\fancyhf{}
\renewcommand{\headrulewidth}{0.5pt}
\renewcommand{\footrulewidth}{0.5pt}
\fancyhead[L]{\small\bfseries\color{teal}\faBook\ SỔ TAY KIẾN THỨC TOÁN 3}
\fancyhead[R]{\small\bfseries\color{teal}Lớp Toán Thầy Lĩnh -- loptoanthaylinh.vn}
\fancyfoot[C]{\small\color{gray}-- Sổ tay kiến thức: Bài 28. Bài toán giải bằng hai bước tính --}

\definecolor{primary}{HTML}{005F73}
\definecolor{secondary}{HTML}{0A9396}
\definecolor{accent}{HTML}{EE9B00}
\definecolor{danger}{HTML}{AE2012}
\definecolor{boxbg}{HTML}{F4F6F8}

\begin{document}

% --- BANNER TIÊU ĐỀ ---
\begin{tcolorbox}[
  enhanced,
  colback=white,
  colframe=primary,
  arc=3mm,
  boxrule=1.2pt,
  left=4mm, right=4mm, top=2.5mm, bottom=2.5mm,
  drop shadow=black!10!white
]
  \centering
  {\Large\bfseries\color{primary}\faBookmark\ BÀI 28: BÀI TOÁN GIẢI BẰNG HAI BƯỚC TÍNH}\\[1.5mm]
  {\small\itshape\color{black!75}Kĩ năng phân tích ngược, tìm ẩn số trung gian và trình bày lời giải 2 bước mạch lạc trên 1 trang A4}
\end{tcolorbox}

\vspace{-2mm}


\noindent
\begin{minipage}[t]{0.49\textwidth}\vspace{0pt}
\begin{tcolorbox}[
  enhanced,
  colback=blue!4!white,
  colframe=blue!75!black,
  arc=2mm, boxrule=0.8pt,
  left=2mm, right=2mm, top=1.8mm, bottom=1.8mm,
  title={\footnotesize\bfseries\color{white}\faKey\ 1. Quy trình 3 bước giải bài toán 2 bước}
]
{\footnotesize
\begin{itemize}\setlength{\itemsep}{1.5pt}\setlength{\leftskip}{-2.5mm}
  \item \textbf{Bước 1: Tìm hiểu \& Tóm tắt đề bài:}
  Xác định đại lượng đã biết và đại lượng câu hỏi yêu cầu tìm.
  \item \textbf{Bước 2: Tìm đại lượng trung gian (Bước 1):}
  Đại lượng chưa cho trực tiếp nhưng cần có để trả lời câu hỏi chính.
  \item \textbf{Bước 3: Trả lời câu hỏi chính (Bước 2):}
  Kết hợp đại lượng ban đầu và đại lượng trung gian để tìm đáp số.
\end{itemize}
}
\end{tcolorbox}
\end{minipage}
\hfill
\begin{minipage}[t]{0.49\textwidth}\vspace{0pt}
\begin{tcolorbox}[
  enhanced,
  colback=teal!4!white,
  colframe=teal!75!black,
  arc=2mm, boxrule=0.8pt,
  left=2mm, right=2mm, top=1.8mm, bottom=1.8mm,
  title={\footnotesize\bfseries\color{white}\faSliders\ 2. Các dạng toán 2 bước điển hình}
]
{\footnotesize
\begin{itemize}\setlength{\itemsep}{1.5pt}\setlength{\leftskip}{-2.5mm}
  \item \textbf{Dạng 1 (Nhiều hơn / Ít hơn rồi tính tổng):}
  \begin{itemize}\setlength{\itemsep}{0.5pt}\setlength{\leftskip}{-2.5mm}
    \item Bước 1: Tìm phần tử thứ hai ($+$ hoặc $-$).
    \item Bước 2: Tính tổng cả hai phần tử ($+$).
  \end{itemize}
  \item \textbf{Dạng 2 (Gấp / Giảm số lần rồi tính tổng/hiệu):}
  \begin{itemize}\setlength{\itemsep}{0.5pt}\setlength{\leftskip}{-2.5mm}
    \item Bước 1: Tìm phần tử thứ hai ($\times$ hoặc $:$).
    \item Bước 2: Tính tổng hoặc còn lại ($+$ hoặc $-$).
  \end{itemize}
\end{itemize}
}
\end{tcolorbox}
\end{minipage}

\vspace{-2.5mm}
% --- KHỐI 3: SƠ ĐỒ TƯ DUY TIKZ ---
\noindent
\begin{tcolorbox}[
  enhanced,
  colback=orange!4!white,
  colframe=orange!80!black,
  arc=2mm, boxrule=0.8pt,
  left=2mm, right=2mm, top=1.5mm, bottom=1.5mm,
  title={\footnotesize\bfseries\color{white}\faSitemap\ SƠ ĐỒ TƯ DUY: CẤU TRÚC PHÂN TÍCH BÀI TOÁN 2 BƯỚC}
]
\centering
\begin{tikzpicture}[
  box/.style={draw, rounded corners=2pt, align=center, font=\scriptsize\bfseries, inner sep=2.5pt, minimum height=0.75cm},
  arr/.style={-{Stealth[length=2mm]}, thick}
]
  \node[box, fill=purple!15, draw=purple!80!black, text width=3.2cm] (s1) at (0, 0) {ĐỀ BÀI CHO\\Đại lượng đã biết $A$};
  \node[box, fill=teal!15, draw=teal!80!black, text width=3.4cm] (s2) at (4.5, 0) {BƯỚC 1: TÌM TRUNG GIAN\\Đại lượng $B$\\($B = A \times k$ hoặc $A : k$ hoặc $A \pm k$)};
  \node[box, fill=orange!15, draw=orange!80!black, text width=3.4cm] (s3) at (9.2, 0) {BƯỚC 2: CÂU HỎI CHÍNH\\Tính Tổng $(A + B)$\\hoặc Hiệu $(A - B)$};

  \draw[arr, purple!80!black] (s1) -- (s2);
  \draw[arr, teal!80!black] (s2) -- (s3);
\end{tikzpicture}
\end{tcolorbox}

\vspace{-2.5mm}
% --- KHỐI 4: VÍ DỤ & MẸO TRÁNH BẪY ---
\noindent
\begin{minipage}[t]{0.49\textwidth}\vspace{0pt}
\begin{tcolorbox}[
  enhanced,
  colback=green!4!white,
  colframe=green!60!black,
  arc=2mm, boxrule=0.8pt,
  left=2mm, right=2mm, top=1.8mm, bottom=1.8mm,
  title={\footnotesize\bfseries\color{white}\faCheckCircle\ Ví dụ thực hành chuẩn mực}
]
{\footnotesize
\textbf{Bài toán mẫu:}\\
Ngăn trên có $12$ quyển sách. Ngăn dưới có số sách gấp $3$ lần ngăn trên. Hỏi cả hai ngăn có bao nhiêu quyển sách?\\
\textbf{Bài giải chuẩn:}\\
Ngăn dưới có số quyển sách là:
$$12 \times 3 = 36\text{ (quyển)}$$
Cả hai ngăn có số quyển sách là:
$$12 + 36 = \mathbf{48}\text{ (quyển)}$$
\phantom{Bài giải: }\textbf{Đáp số:} $48$ quyển sách.
}
\end{tcolorbox}
\end{minipage}
\hfill
\begin{minipage}[t]{0.49\textwidth}\vspace{0pt}
\begin{tcolorbox}[
  enhanced,
  colback=red!4!white,
  colframe=red!70!black,
  arc=2mm, boxrule=0.8pt,
  left=2mm, right=2mm, top=1.8mm, bottom=1.8mm,
  title={\footnotesize\bfseries\color{white}\faExclamationTriangle\ Mẹo tránh bẫy trắc nghiệm}
]
{\footnotesize
\begin{itemize}\setlength{\itemsep}{1.5pt}\setlength{\leftskip}{-2.5mm}
  \item \textbf{Bẫy dừng lại ở bước 1 (Bẫy phổ biến nhất):}\\
  Khi tính xong đại lượng trung gian (ví dụ tính ra ngăn dưới có $36$ quyển), học sinh vội vàng lấy đó làm đáp số mà quên mất câu hỏi là \textbf{"cả hai ngăn"}.
  \item \textbf{Bí quyết kiểm tra:}\\
  Trước khi ghi đáp số, hãy đọc lại câu hỏi của đề bài một lần nữa: "Đề hỏi một đối tượng hay hỏi cả hai?".
\end{itemize}
}
\end{tcolorbox}
\end{minipage}


\end{document}
